%% kapa: exam=1 round=32 year=2021 period=1 form=A subject=economics subject_ko=경제학원론 seq=23 number=63
%% source: Q-Net 감정평가사 기출문제 (cst00309) artlSeq=5208397 file=extracted/1008/5208397/A형/2021년도 제32회 감정평가사 1차 1교시 A형.pdf
%% conversion: kordoc 4.24.51 -> tools/gen1.py (auto); review=proofread-sonnet; flags=none
\begin{question}{63}
화폐에 관한 설명으로 옳은 것은?
\begin{choices}
  \item 상품화폐의 내재적 가치는 변동하지 않는다.
  \item M2는 준화폐(near money)를 포함하지 않는다.
  \item 명령화폐(fiat money)는 내재적 가치를 갖는 화폐이다.
  \item 가치 저장수단의 역할로 소득과 지출의 발생 시점을 분리시켜 준다.
  \item 다른 용도로 사용될 수 있는 재화는 교환의 매개 수단으로 활용될 수 없다.
\end{choices}
\end{question}
\answer{4}
